% ================= ROZDZIAŁ 4.3 =================
\subsection{Wymagania integracyjne}

Wymagania integracyjne opisują sposób współpracy systemu \texttt{WorkFlow} z usługami technicznymi, komponentami infrastruktury oraz zewnętrznymi źródłami zdarzeń. W aktualnym zakresie projektu integracje dotyczą przede wszystkim bazy danych, magazynu plików, czytników RCP, komunikacji pomiędzy frontendem i backendem oraz środowiska uruchomieniowego opartego o kontenery.

Celem wymagań integracyjnych jest określenie, jakie elementy muszą ze sobą współpracować, aby system działał jako kompletna aplikacja webowa. Wymagania te nie opisują pojedynczych funkcji użytkownika, lecz zależności pomiędzy usługami i warstwami systemu.

\subsubsection{Integracja z bazą danych}

System powinien integrować się z relacyjną bazą danych PostgreSQL. Baza danych stanowi główne miejsce przechowywania informacji biznesowych, takich jak pracownicy, zakłady, projekty, umowy, certyfikaty, nieobecności, dokumenty, zdarzenia czasu pracy, wpisy czasu pracy oraz eksporty rozliczeniowe.

Komunikacja backendu z bazą danych odbywa się z wykorzystaniem narzędzia Prisma ORM. Schemat bazy został zdefiniowany w pliku \texttt{schema.prisma}, co umożliwia utrzymanie spójności pomiędzy modelem danych aplikacji a strukturą fizycznej bazy danych.

Wymagania w tym obszarze obejmują:
\begin{itemize}
    \item możliwość połączenia backendu z bazą PostgreSQL za pomocą zmiennej środowiskowej,
    \item odwzorowanie encji biznesowych w modelach Prisma,
    \item obsługę relacji pomiędzy tabelami,
    \item możliwość synchronizacji struktury bazy danych ze schematem Prisma,
    \item możliwość inicjalizacji bazy danymi startowymi.
\end{itemize}

\subsubsection{Integracja z magazynem plików}

System powinien współpracować z magazynem obiektowym MinIO. Pliki dokumentów nie są przechowywane bezpośrednio w tabelach bazy danych. W bazie zapisywane są jedynie metadane dokumentu, takie jak nazwa pliku, właściciel dokumentu oraz ścieżka do pliku w magazynie.

Takie rozwiązanie pozwala oddzielić dane relacyjne od plików binarnych. Jest to szczególnie istotne w przypadku dokumentów kadrowych, skanów umów, zaświadczeń, certyfikatów lub dokumentów powiązanych z nieobecnościami.

Wymagania w zakresie integracji z magazynem plików obejmują:
\begin{itemize}
    \item konfigurację endpointu MinIO za pomocą zmiennych środowiskowych,
    \item obsługę danych dostępowych do magazynu plików,
    \item wskazanie nazwy bucketu przechowującego dokumenty,
    \item zapisywanie w bazie danych ścieżki lub adresu pliku,
    \item możliwość powiązania dokumentu z pracownikiem, umową, certyfikatem lub nieobecnością.
\end{itemize}

\subsubsection{Integracja z czytnikami RCP}

System powinien udostępniać techniczny endpoint przeznaczony do odbierania zdarzeń z czytników RCP. Czytnik jest powiązany z zakładem oraz projektem, dzięki czemu zdarzenia wejścia i wyjścia mogą zostać przypisane do konkretnego pracownika oraz miejsca pracy.

Zdarzenie przekazywane do systemu powinno zawierać:
\begin{itemize}
    \item numer seryjny czytnika,
    \item identyfikator karty RFID,
    \item typ zdarzenia: \texttt{IN} albo \texttt{OUT},
    \item czas zdarzenia,
    \item opcjonalny zewnętrzny identyfikator zdarzenia.
\end{itemize}

Endpoint urządzeń powinien być zabezpieczony tokenem technicznym, niezależnym od tokenów użytkowników logujących się do aplikacji. Pozwala to oddzielić komunikację maszynową od standardowej komunikacji użytkownika z systemem.

Po odebraniu zdarzenia backend powinien:
\begin{enumerate}
    \item zweryfikować token urządzenia,
    \item odnaleźć aktywny czytnik na podstawie numeru seryjnego,
    \item odnaleźć aktywnego pracownika na podstawie identyfikatora karty RFID,
    \item zapisać surowe zdarzenie w tabeli \texttt{TimeEvent},
    \item w przypadku zdarzenia \texttt{OUT} spróbować utworzyć wpis czasu pracy na podstawie wcześniejszego zdarzenia \texttt{IN}.
\end{enumerate}

\subsubsection{Komunikacja pomiędzy frontendem i backendem}

Frontend powinien komunikować się z backendem za pomocą API HTTP. Backend udostępnia punkty końcowe odpowiedzialne za logowanie, pobieranie danych, tworzenie rekordów, aktualizację informacji oraz obsługę operacji biznesowych.

Frontend odpowiada za:
\begin{itemize}
    \item prezentację danych użytkownikowi,
    \item obsługę formularzy,
    \item wyświetlanie list i szczegółów rekordów,
    \item kierowanie użytkownika do odpowiednich widoków,
    \item przekazywanie danych do backendu.
\end{itemize}

Backend odpowiada za:
\begin{itemize}
    \item autoryzację żądań,
    \item walidację danych wejściowych,
    \item realizację logiki biznesowej,
    \item komunikację z bazą danych,
    \item komunikację z magazynem plików,
    \item obsługę technicznych endpointów, takich jak odbieranie zdarzeń RCP.
\end{itemize}

Taki podział umożliwia oddzielenie warstwy prezentacji od logiki biznesowej i ułatwia dalsze utrzymanie aplikacji.

\subsubsection{Integracja przez Nginx}

W środowisku Docker Compose system wykorzystuje Nginx jako reverse proxy. Oznacza to, że użytkownik korzysta z jednego adresu aplikacji, a Nginx przekierowuje ruch do odpowiedniej usługi.

Wymagania dotyczące Nginx obejmują:
\begin{itemize}
    \item przekierowanie ruchu aplikacji do frontendu,
    \item przekierowanie zapytań API do backendu,
    \item uproszczenie adresów wykorzystywanych przez użytkownika,
    \item oddzielenie wewnętrznych nazw usług od adresu widocznego z poziomu przeglądarki.
\end{itemize}

Przykładowy podział ruchu wygląda następująco:
\begin{itemize}
    \item ścieżka \texttt{/} kierowana jest do frontendu,
    \item ścieżka \texttt{/api} kierowana jest do backendu.
\end{itemize}

Dzięki temu użytkownik nie musi znać portów poszczególnych usług, a aplikacja może być dostępna pod jednym adresem.

\subsubsection{Środowisko Docker Compose}

System powinien być możliwy do uruchomienia za pomocą Docker Compose. Plik konfiguracyjny opisuje usługi niezbędne do działania aplikacji oraz zależności pomiędzy nimi.

Środowisko powinno obejmować:
\begin{itemize}
    \item usługę frontendu,
    \item usługę backendu,
    \item usługę PostgreSQL,
    \item usługę MinIO,
    \item usługę Nginx.
\end{itemize}

Konfiguracja powinna korzystać ze zmiennych środowiskowych przechowywanych w pliku \texttt{.env}. W repozytorium powinien znajdować się wyłącznie plik przykładowy \texttt{.env.example}, który opisuje wymagane zmienne bez ujawniania rzeczywistych haseł, tokenów i sekretów.

\subsubsection{Inicjalizacja danych}

System powinien umożliwiać przygotowanie bazy danych do pracy po pierwszym uruchomieniu. W projekcie przewidziano dwa mechanizmy inicjalizacji danych:
\begin{itemize}
    \item \textbf{seed minimalny} -- tworzący podstawowy zakład oraz konta startowe dla głównych ról użytkowników,
    \item \textbf{seed demonstracyjny} -- tworzący komplet danych potrzebnych do prezentacji działania systemu.
\end{itemize}

Seed minimalny jest przeznaczony do czystego uruchomienia aplikacji. Seed demonstracyjny służy natomiast do przygotowania środowiska pokazowego i zawiera między innymi pracowników, projekty, czytniki RCP, umowy, certyfikaty, nieobecności, dokumenty, wpisy czasu pracy oraz przykładowy eksport rozliczeniowy.

Ze względu na to, że seed demonstracyjny może czyścić istniejące dane, jego uruchomienie powinno wymagać dodatkowej zmiennej środowiskowej potwierdzającej intencję wykonania operacji.

\subsubsection{Integracje możliwe do rozbudowy}

Aktualna wersja systemu skupia się na podstawowych integracjach wymaganych do działania aplikacji. Przyjęta architektura pozwala jednak na późniejsze rozszerzenie systemu o kolejne mechanizmy.

Do możliwych kierunków rozwoju należą:
\begin{itemize}
    \item integracja z zewnętrznym systemem księgowym,
    \item wysyłka powiadomień e-mail,
    \item integracja z zewnętrznym dostawcą tożsamości,
    \item rozbudowa harmonogramów zadań cyklicznych,
    \item bardziej zaawansowany eksport danych rozliczeniowych.
\end{itemize}

Wymienione elementy nie są wymagane do działania aktualnej wersji projektu. Zostały wskazane jako potencjalne rozszerzenia, które mogą zostać dodane w przyszłości bez zmiany podstawowych założeń architektonicznych systemu.